%! Characteristics of Linear Functions — Unit 2, Lesson 2.0 — Group Activity (Answer Key)
\documentclass[10pt]{article}
\usepackage{algebra2-article}
\usepackage{algebra2-key}   % pulls in -boxes; do NOT also load -boxes
\newcommand{\TallMath}[1]{$\displaystyle #1\rule[-1.4em]{0pt}{3.2em}$}
\newcommand{\lgrid}[4]{%
  \draw[step=1, linegray!40, very thin] (#1,#3) grid (#2,#4);
  \draw[->, semithick, charcoal] (#1,0)--(#2,0) node[below right, font=\scriptsize]{$x$};
  \draw[->, semithick, charcoal] (0,#3)--(0,#4) node[left, font=\scriptsize]{$y$};}

\begin{document}

\pageheader{Unit 2, Lesson 2.0}{Group Activity --- Answer Key}
\namepartnerperiod

\vspace{0.08in}

\begin{headlinebox}{forestbg}
\small\textbf{Reading Lines.} Every characteristic of a linear function is something you can
\emph{point to} on its graph. With your partner, read each graph carefully and \emph{justify} your
answers --- especially where a function is positive or negative.
\end{headlinebox}

\vspace{0.06in}

\begin{tcolorbox}[colback=white, colframe=black!40, title=\textbf{Tier R --- Remediate},
    fonttitle=\bfseries, arc=1.5mm, left=3mm, right=3mm, top=2mm, bottom=2mm, breakable]
The graph shows the line $\ell(x) = x - 3$.
\begin{center}
\begin{tikzpicture}[scale=0.5]
  \lgrid{-2}{5}{-5}{3}
  \draw[very thick, forest] (-1,-4)--(5,2);
  \fill[charcoal] (0,-3) circle (3pt);
  \fill[charcoal] (3,0) circle (3pt);
\end{tikzpicture}
\end{center}
\begin{enumerate}[leftmargin=1.7em, itemsep=7pt]
  \item The line crosses the $y$-axis at $(0,\ \textcolor{keyred}{\mathbf{-3}})$ and the $x$-axis at
        $(\textcolor{keyred}{\mathbf{3}},\ 0)$.
  \item Starting at $(3,0)$, go \emph{right} $1$ and count how far \emph{up} you go to return to the
        line. Slope $= \dfrac{\text{rise}}{\text{run}} = \textcolor{keyred}{\mathbf{1}}$.
  \item Reading left to right, does the line \textbf{rise} or \textbf{fall}? \ans{rises} \quad
        So it is (circle one): \textcolor{keyred}{\textbf{increasing}} \; / \; \textbf{decreasing}.
\end{enumerate}
\end{tcolorbox}

\begin{tcolorbox}[colback=white, colframe=black!40, title=\textbf{Tier A --- Approaching Proficiency},
    fonttitle=\bfseries, arc=1.5mm, left=3mm, right=3mm, top=2mm, bottom=2mm, breakable]
For \emph{each} line, give the full feature list and \emph{justify} the positive/negative intervals.
\begin{center}
\begin{tikzpicture}[scale=0.44]
  % Line 1: y = 2x + 2
  \begin{scope}
    \lgrid{-4}{4}{-4}{6}
    \draw[very thick, forest] (-3,-4)--(2,6);
    \fill[charcoal] (0,2) circle (3pt);
    \fill[charcoal] (-1,0) circle (3pt);
    \node[charcoal, font=\scriptsize] at (0,-5.2){Line 1: $f(x) = 2x + 2$};
  \end{scope}
  % Line 2: y = -1/2 x + 2
  \begin{scope}[xshift=10cm]
    \lgrid{-3}{6}{-3}{5}
    \draw[very thick, forest] (-2,3)--(6,-1);
    \fill[charcoal] (0,2) circle (3pt);
    \fill[charcoal] (4,0) circle (3pt);
    \node[charcoal, font=\scriptsize] at (1.5,-4.2){Line 2: $g(x) = -\tfrac{1}{2}x + 2$};
  \end{scope}
\end{tikzpicture}
\end{center}
\renewcommand{\arraystretch}{1.5}
\begin{tabularx}{\linewidth}{@{}>{\bfseries}p{0.30\linewidth} X X@{}}
 & \textbf{Line 1: $f(x)=2x+2$} & \textbf{Line 2: $g(x)=-\tfrac12 x+2$} \\
\hline
$y$-intercept & \ans{$(0,2)$} & \ans{$(0,2)$} \\
$x$-intercept (zero) & \ans{$(-1,0)$} & \ans{$(4,0)$} \\
Slope & \ans{$2$} & \ans{$-\tfrac12$} \\
Increasing / decreasing & \ans{increasing} & \ans{decreasing} \\
Positive when & \ans{$x>-1$} & \ans{$x<4$} \\
Negative when & \ans{$x<-1$} & \ans{$x>4$} \\
\end{tabularx}
\vspace{2pt}

\textbf{Justify:} For Line~2, explain why $g$ is \emph{positive} to the left of its zero but the
line is still \emph{decreasing}. (These are different ideas!)
\ansline{``Positive'' asks where the graph sits: above the $x$-axis, which is $x<4$ (left of the zero). ``Decreasing'' asks which way it moves: as $x$ increases the output falls (slope $-\tfrac12<0$). So $g$ can be above the axis \emph{and} heading downward at the same time --- sign of the output vs.\ direction of travel are separate questions.}
\end{tcolorbox}

\begin{tcolorbox}[colback=white, colframe=black!40, title=\textbf{Tier E --- Extension},
    fonttitle=\bfseries, arc=1.5mm, left=3mm, right=3mm, top=2mm, bottom=2mm, breakable]
\textbf{Back to the tank.} The volume of a draining tank is $V(t) = -3t + 12$ gallons after $t$
minutes.
\begin{center}
\begin{tikzpicture}[scale=0.5]
  \draw[step=1, linegray!40, very thin] (0,0) grid (5,13);
  \draw[->, semithick, charcoal] (0,0)--(5,0) node[below right, font=\scriptsize]{$t$ (min)};
  \draw[->, semithick, charcoal] (0,0)--(0,13) node[left, font=\scriptsize]{$V$ (gal)};
  \foreach \y in {3,6,9,12} \node[charcoal, font=\scriptsize, left] at (0,\y){\y};
  \foreach \x in {1,2,3,4} \node[charcoal, font=\scriptsize, below] at (\x,0){\x};
  \draw[very thick, forest] (0,12)--(4,0);
  \fill[charcoal] (0,12) circle (3pt);
  \fill[charcoal] (4,0) circle (3pt);
\end{tikzpicture}
\end{center}
\begin{enumerate}[leftmargin=1.7em, itemsep=8pt]
  \item Interpret each feature \emph{in context}: the $y$-intercept means \ans{the tank started with $12$ gallons};
        the slope means \ans{it loses $3$ gallons each minute}; the $x$-intercept (zero) means \ans{the tank is empty at $t=4$ min}.
  \item What is the \emph{meaningful domain} of $V$ in this situation, and why is it \emph{not} all
        real numbers?
        \ansline{$0\le t\le 4$. Time before $t=0$ isn't measured, and after $t=4$ the formula gives a negative volume, which is impossible --- so only $0$ to $4$ minutes make sense.}
  \item \textbf{Horizontal special case.} A second, sealed tank never drains: its volume is the
        constant function $C(t) = 8$. What is its slope? Is it increasing, decreasing, or constant?
        What is its range, and does it have a zero?
        \ansline{Slope $0$; constant (neither increasing nor decreasing); range $\{8\}$ (only the value $8$); no zero --- it never reaches $V=0$.}
\end{enumerate}
\end{tcolorbox}

\begin{teachernote}
Tier~A's justify item is the crux of the lesson: keep pushing groups to separate ``sign of the
output'' (positive/negative) from ``direction of travel'' (increasing/decreasing). Line~2 makes the
point on its own --- decreasing yet positive on $(-\infty,4)$. In Tier~E, item~2 (meaningful domain)
is where the modeling payoff lives; accept any clear reason that negative time or negative volume is
impossible.
\end{teachernote}

\end{document}
